% =========================================================================
% Chapter 4: Results
% ZOI digits locked to the laboratory file. C1 = 50 mg/mL, C2 = 100 mg/mL.
% =========================================================================

\chapter{Results}\label{ch:results}

\section{Extract yield and appearance}

Soxhlet extraction of \SI{50}{\gram} of powdered aerial material with \SI{500}{\milli\litre} of distilled water gave a dark, concentrated aqueous extract corresponding to 24.13\% (w/w) of the starting powder. Working solutions C1 (\SI{50}{\milli\gram\per\milli\litre}) and C2 (\SI{100}{\milli\gram\per\milli\litre}) were brown to black.

\section{Qualitative phytochemical profile}

All nine labelled tubes on the rack gave a positive reading (\cref{tab:phytochem}, \cref{fig:phytochem}). Flavonoids were yellow-brown after alkaline reagent. Terpenoids and quinones were dark. Alkaloids gave a dull deposit with Mayer's reagent. Tannins and phenols were blue-black with ferric chloride. Saponins held a froth in the tube. Carbohydrates and phytosterols were also scored positive. No class on the rack was negative.

\begin{table}[htbp]
\centering
\caption{Qualitative phytochemical analysis of the aqueous extract of \taxa{Dipsacus inermis}.}
\label{tab:phytochem}
\small
\begin{tabular}{@{}>{\raggedright\arraybackslash}p{0.18\textwidth}
  >{\raggedright\arraybackslash}p{0.16\textwidth}
  >{\raggedright\arraybackslash}p{0.21\textwidth}
  >{\raggedright\arraybackslash}p{0.19\textwidth}
  >{\raggedright\arraybackslash}p{0.14\textwidth}@{}}
\toprule
\textbf{Class} & \textbf{Test} & \textbf{Reagents} & \textbf{Observation} & \textbf{Inference} \\
\midrule
Alkaloids & Mayer's & Mayer's reagent & Dull white precipitate & Positive (+) \\
Flavonoids & Alkaline reagent & 10\% NaOH & Yellow colour & Positive (+) \\
Phenols & Ferric chloride & 10\% FeCl$_3$ & Bluish black & Positive (+) \\
Tannins & Ferric chloride & FeCl$_3$ & Dark blue / greenish black & Positive (+) \\
Saponins & Froth & Water, shaken & Persistent froth & Positive (+) \\
Terpenoids & Salkowski & Thionyl chloride (notebook) & Pink / dark colour & Positive (+) \\
Phytosterols & Salkowski & Chloroform + H$_2$SO$_4$ & Reddish brown & Positive (+) \\
Carbohydrates & Fehling's & Fehling A and B, heated & Brick-red precipitate & Positive (+) \\
Quinones & Sulphuric acid & H$_2$SO$_4$ & Red / dark colour & Positive (+) \\
\bottomrule
\end{tabular}
\end{table}

\begin{figure}[htbp]
\centering
\includegraphics[width=0.95\textwidth]{fig8_phytochem_tubes.jpg}
\caption{Qualitative phytochemical reactions of the aqueous extract of \taxa{Dipsacus inermis}. Tubes, left to right: flavonoids, terpenoids, quinones, alkaloids, tannins, saponins (froth visible), phenols, carbohydrates, phytosterols.}
\label{fig:phytochem}
\end{figure}

\section{Zone of inhibition}

Distilled water (N) gave $0.00 \pm 0.00$\,mm against all three isolates. Both extract concentrations produced measurable zones (\cref{tab:zoi}). Zones were measured inclusive of the disc. Mean values are from the triplicate readings and are not rounded beyond two decimal places.

Against \taxa{Escherichia coli} the C1 mean was $11.94 \pm 0.14$\,mm and the C2 mean was $14.99 \pm 0.51$\,mm. The GEN 5 disc measured $17.00 \pm 0.00$\,mm.

Against \taxa{Pseudomonas aeruginosa} the C1 mean was $16.03 \pm 0.25$\,mm and the C2 mean was $12.53 \pm 0.45$\,mm. The GEN 5 disc measured $17.17 \pm 0.29$\,mm. C1 was the largest extract zone in the study. C2 was smaller than C1 for this isolate.

Against \taxa{Klebsiella pneumoniae} the C1 mean was $12.35 \pm 0.56$\,mm and the C2 mean was $13.25 \pm 0.43$\,mm. The GEN 5 disc measured $15.17 \pm 0.29$\,mm.

Incubated plates marked N, C1, C2, and GEN 5 are shown in \cref{fig:finalplates}. Species names are not written on those two plates. Assay plates dated 04 August 2026 are shown in \cref{fig:assaytrio}. The upper-left lid reads \textit{K.\ pneumoniae}. The right lid reads Pseudomonas. The species line on the lower lid is not fully legible and is not named here. A representative GEN 5 plate is shown in \cref{fig:gen5}. Mean values with standard-deviation bars from \cref{tab:zoi} are plotted in \cref{fig:zoibars}.

\begin{table}[htbp]
\centering
\caption{Antibacterial activity (zone of inhibition, mm) of the aqueous extract of \taxa{Dipsacus inermis}. Values are mean $\pm$ SD ($n = 3$), measured inclusive of the disc. C1 = \SI{50}{\milli\gram\per\milli\litre}; C2 = \SI{100}{\milli\gram\per\milli\litre}; GEN 5 = gentamicin \SI{5}{\micro\gram} disc; N = distilled water.}
\label{tab:zoi}
\small
\begin{tabular}{@{}l*{4}{S[table-format=2.2(2)]}@{}}
\toprule
Isolate & {N (water)} & {C1} & {C2} & {GEN 5} \\
\midrule
\textit{E.\ coli} & 0.00 \pm 0.00 & 11.94 \pm 0.14 & 14.99 \pm 0.51 & 17.00 \pm 0.00 \\
\textit{P.\ aeruginosa} & 0.00 \pm 0.00 & 16.03 \pm 0.25 & 12.53 \pm 0.45 & 17.17 \pm 0.29 \\
\textit{K.\ pneumoniae} & 0.00 \pm 0.00 & 12.35 \pm 0.56 & 13.25 \pm 0.43 & 15.17 \pm 0.29 \\
\bottomrule
\end{tabular}
\end{table}

\begin{figure}[htbp]
\centering
\includegraphics[width=0.92\textwidth]{fig11_final_result_plates.png}
\caption{Two LB assay plates after incubation, marked distilled water (N), C1, C2, and gentamicin GEN 5. Clearance is visible around the extract wells and the GEN 5 disc. No binomial is written on these plates.}
\label{fig:finalplates}
\end{figure}

\begin{figure}[htbp]
\centering
\includegraphics[width=0.88\textwidth]{fig9_assay_trio_dated.jpg}
\caption{Assay plates, 04 August 2026. The upper-left lid reads \textit{K.\ pneumoniae}. The right lid reads Pseudomonas. The species line on the lower plate is not fully legible and is not assigned here. Positions C1, C2, N, and a paper disc are present on each plate.}
\label{fig:assaytrio}
\end{figure}

\begin{figure}[htbp]
\centering
\includegraphics[width=0.62\textwidth]{fig10_gen5_plate.jpg}
\caption{Representative plate showing gentamicin GEN 5 (P), sterile water (N) with no halo, and two extract positions with surrounding clearance. No binomial is written on this close-up.}
\label{fig:gen5}
\end{figure}

\begin{figure}[htbp]
\centering
\begin{tikzpicture}
\begin{axis}[
  ybar,
  bar width=7pt,
  ylabel={Zone of inhibition (mm)},
  ymin=0,
  ymax=22,
  width={0.98\textwidth},
  height=8.2cm,
  enlarge x limits=0.18,
  symbolic x coords={Ecoli, Paer, Kpneu},
  xtick=data,
  xticklabels={\textit{E.\ coli}, \textit{P.\ aeruginosa}, \textit{K.\ pneumoniae}},
  legend style={at={(0.5,-0.22)}, anchor=north, legend columns=3, font=\small, draw=none},
  grid=major,
  grid style={dashed, gray!30},
  /pgf/number format/fixed,
  error bars/y dir=both,
  error bars/y explicit,
]
\addplot+[fill=dissertationblue!70, draw=black, error bars/.cd, y dir=both, y explicit]
  coordinates {(Ecoli,11.94) +-(0.14,0.14) (Paer,16.03) +-(0.25,0.25) (Kpneu,12.35) +-(0.56,0.56)};
\addplot+[fill=dissertationgreen!70, draw=black, error bars/.cd, y dir=both, y explicit]
  coordinates {(Ecoli,14.99) +-(0.51,0.51) (Paer,12.53) +-(0.45,0.45) (Kpneu,13.25) +-(0.43,0.43)};
\addplot+[fill=dissertationgray!50, draw=black, error bars/.cd, y dir=both, y explicit]
  coordinates {(Ecoli,17.00) +-(0.00,0.00) (Paer,17.17) +-(0.29,0.29) (Kpneu,15.17) +-(0.29,0.29)};
\legend{C1 (50 mg/mL), C2 (100 mg/mL), GEN 5}
\end{axis}
\end{tikzpicture}
\caption{Clustered bar graph of mean zone of inhibition ($\pm$ SD) for C1 (\SI{50}{\milli\gram\per\milli\litre}), C2 (\SI{100}{\milli\gram\per\milli\litre}), and gentamicin GEN 5 against the three isolates.}
\label{fig:zoibars}
\end{figure}

Raw triplicate readings are listed in \cref{app:zoi}.

\FloatBarrier
